\begin{frame}{Background: Parameter-Efficient Fine-Tuning (PEFT)}
\begin{itemize}
  \item Large language models achieve strong transfer performance, but full fine-tuning updates all model parameters.
  \item Low-Rank Adaptation (LoRA) reduces trainable parameters by inserting low-rank update matrices:
  \[
    W' = W + \Delta W,\quad \Delta W = BA,\quad
    A \in \mathbb{R}^{r \times d},\ B \in \mathbb{R}^{k \times r}.
  \]
  \item The rank \(r\) controls the adaptation capacity and computational cost.
  \item In low-resource settings, a fixed rank can be too large for simple or noisy datasets and too small for complex datasets.
\end{itemize}
\end{frame}

\begin{frame}{Problem Setting}
\begin{block}{Research Question}
Can AdaLoRA be improved by selecting adapter capacity from dataset characteristics before fine-tuning?
\end{block}

\vspace{8pt}
\begin{itemize}
  \item Standard LoRA: uses a fixed rank and fixed target modules.
  \item Standard AdaLoRA: adaptively reallocates rank during training, but the initial configuration is usually hand-designed.
  \item Proposed method: Star-LoRA estimates dataset difficulty before training and uses it to configure AdaLoRA.
  \item Final task in this project: fake review detection with ModernBERT-large.
\end{itemize}
\end{frame}

\begin{frame}{Motivation}
\begin{itemize}
  \item Low-resource fine-tuning is sensitive to overfitting and unstable gradient signals.
  \item Different datasets have different properties:
  \begin{itemize}
    \item sample count,
    \item label imbalance,
    \item token sparsity,
    \item token-frequency skew,
    \item label entropy,
    \item input-embedding variance.
  \end{itemize}
  \item These statistics can indicate whether the model needs smaller, larger, or broader adapter capacity.
\end{itemize}
\end{frame}
